\section{Le moteur, et ce qu'il a fait sur le yen}

\subsection{La règle est celle de l'or, l'instrument ne l'est pas}

Le moteur appliqué ici est \emph{exactement} celui de la sleeve or~: un
momentum de série temporelle qui achète quand le passé récent monte et sort
quand il cesse de monter. Il moyenne le \emph{signe} de quatre horizons fixes
--- 40, 60, 120 et 250~séances --- plutôt que d'en choisir un, et la position
s'ouvre quand le score devient positif, se ferme quand il redevient négatif. La
taille est pilotée par la volatilité, avec la même cible de 55\,\% et le même
plafond de 6,6$\times$.

Rien n'a été réglé pour le yen. C'est délibéré~: la seule question posée à ce
candidat était de savoir si le moteur \emph{existant} tient sur un autre
sous-jacent, et tout réglage spécifique aurait transformé une vérification en
recherche de paramètres. Le seul changement est l'entrée
\code{Inp\_Gold\_Symbol=USDJPY}, qui pointe le même code sur un autre
instrument~; c'est pourquoi les journaux du simulateur portent encore
l'étiquette \code{Gold\_Momentum}.

\begin{table}[H]
\centering
\begin{tabular}{@{}lll@{}}
\toprule
Paramètre & Valeur & Rôle \\
\midrule
Instrument        & USDJPY.c                     & CFD de change du courtier \\
Horizons          & 40 / 60 / 120 / 250 séances  & moyennés, jamais sélectionnés \\
Vente à découvert & désactivée                   & la jambe short a été tuée au cycle \\
Cible de volatilité & 55\,\%                     & échelle de position \\
Levier maximal    & 6,6$\times$                  & la contrainte effective ici \\
Stop de sécurité  & 4\,\%                        & jamais déclenché sur ce run \\
Glissement        & 2 points de base par côté    & hypothèse d'exécution \\
Allocation        & 100\,\% (sleeve isolée)      & facteur de risque 1,0 \\
Dépôt             & 100\,000\,\$                 & sous-compte de recherche \\
Borne de séance   & 21{:}00 UTC                  & clôture de séance du courtier \\
\bottomrule
\end{tabular}
\caption{Configuration du run de recherche USD/JPY --- ce n'est pas une
configuration de production}
\end{table}

\subsection{Les trente-cinq trades dans leur contexte}

La figure~\ref{fig:usdjpy-timeline} montre tout le run d'un coup~: le cours
d'USD/JPY, les périodes en position, le score de décision et le levier
appliqué.

\begin{figure}[H]
\centering
\includegraphics[width=\linewidth]{usdjpy_trades_timeline}
\caption{Les 35 trades dans leur contexte de prix. Bandes vertes~: positions
gagnantes~; bandes rouges~: perdantes. Panneau du milieu~: le score d'ensemble
--- il passe par $\pm$0,25 et $\pm$0,5 parce qu'il moyenne quatre signes.
Panneau du bas~: le levier calculé chaque séance, et sa valeur retenue à
l'entrée de chaque trade.}
\label{fig:usdjpy-timeline}
\end{figure}

Trois lectures se dégagent immédiatement.

D'abord, \textbf{la fenêtre s'ouvre sur huit mois de silence}. La première
position ne tombe que le 14~septembre~2021, alors que la simulation démarre le
1\textsuperscript{er}~janvier. Ce n'est pas une décision du moteur mais une
contrainte d'historique~: il lui faut 252~barres quotidiennes avant de calculer
son horizon le plus long, et le journal l'annonce dès l'initialisation
(«~34/252 D1 bars~; the sleeve stays flat until history accumulates~»).

Ensuite, \textbf{tout le résultat tient dans la première bande verte}. Le
mégatrend du yen de 2021--2022 --- de 110 à 140~yens par dollar --- est capté
par une seule position, ouverte dix jours après le premier signal et tenue
jusqu'en novembre~2022. Après elle, trente-quatre trades s'enchaînent sur
trois ans et demi pour un résultat cumulé \emph{négatif}.

Enfin, \textbf{le levier reste collé à son plafond}. Le panneau du bas ne
descend jamais sous 3$\times$ et passe la moitié du temps à 6,6$\times$~: le
change majeur est un actif peu volatil, et un mécanisme calibré pour viser
55\,\% de volatilité annuelle y demande en permanence plus de levier qu'il n'en
obtient. Le chapitre~\ref{sec:sizing} y revient.

\begin{insightbox}
Le score n'est presque jamais à son maximum~: \textbf{33~entrées sur~35} se
font avec un score de \textbf{0,5} --- deux des quatre horizons positifs
seulement --- et deux seulement avec les quatre horizons alignés. Comme sur
l'or, le moteur n'attend pas la certitude, il prend position sur une majorité
simple. Le score du yen oscille beaucoup~: la sleeve n'est en position que
\metric{67\,\%} du temps --- 1\,133~jours sur les 1\,688 que couvre la série de
trades --- et vingt-deux de ses trente-cinq allers-retours durent une semaine
ou moins.
\end{insightbox}
